\documentclass[11pt]{amsart}
\usepackage[T1]{fontenc}
\usepackage[utf8]{inputenc}
\usepackage{amssymb,booktabs}
\usepackage[margin=1in]{geometry}
\usepackage[hidelinks]{hyperref}
\newcommand{\code}[1]{\nolinkurl{#1}}
\newcommand{\Hom}{\mathsf{Hom}}
\newcommand{\isContr}{\mathsf{isContr}}
\setlength{\emergencystretch}{2em}
\AddToHook{shipout/foreground}{%
 \put(0.5\paperwidth,-\paperheight+0.55cm){\makebox[0pt]{\tiny\ttfamily
 unit/doc/hcat-segal-refactor.tex --- 2026-11-01}}}
\title{Homology from a finite semi-Segal presentation}
\date{1 November 2026}
\begin{document}
\maketitle

The experimental root \code{unit/sag/hcat-segal/All.seg} constructs new
composition and associativity operations from horn fillers, proves their
pentagon, and checks the existing homology theorem files against this
interface. The original \code{hcat} development is unchanged.
This establishes the proposed refactor at the level of an internal
category supplied with finite semi-Segal data. Constructing these data
from the original ambient categorical axioms remains a separate task.

\section{The inputs and the selected operations}

For an object type $O$ and an edge family $E:O\to O\to\mathcal U_l$,
take a triangle family
\[
 A_2(f,g,h):\mathcal U_l,\qquad
 \isContr\left(\sum_{h:E(a,c)}A_2(f,g,h)\right)
 \quad(f:E(a,b),\ g:E(b,c)).
\]
Write $fg$ for composition in the order of traversal. Define $fg$ and
its triangle to be the two components of the selected horn filler.
The fundamental theorem of identity types gives, for every $h$,
\[
 A_2(f,g,h)\simeq(fg=h).
\]
This equivalence justifies using composition paths as the triangle
family at the next stage. The normalized horn fillers are singletons.

A $\Lambda^3_1$ horn now consists of four vertices, six edges, and
three paths $q_{012}$, $q_{123}$, $q_{013}$. Its missing triangle is
\[
 q_{023}:e_{02}e_{23}=e_{03}.
\]
Take a family $A_3(u,q_{023})$ over these full boundaries, with
contractible paired horn fillers. On the canonical horn of $f,g,h$,
put
\[
 e_{02}=fg,\quad e_{13}=gh,\quad e_{03}=f(gh),\qquad
 q_{012}=q_{123}=q_{013}=\mathsf{refl}.
\]
The missing triangle selected by the horn condition is
\[
 \alpha(f,g,h):(fg)h=f(gh).
\]
This is \code{Associativity.hcomp-assoc}. No associator is an input
parameter, and no comparison with the former associator is used.

\section{Normalized tetrahedra}

For a complete triangular boundary define
\begin{align*}
 L&=\mathsf{ap}_{(-e_{23})}(q_{012})\cdot q_{023},\\
 R&=\bigl(\alpha(e_{01},e_{12},e_{23})\cdot
           \mathsf{ap}_{(e_{01}-)}(q_{123})\bigr)\cdot q_{013}.
\end{align*}
Both paths go from $(e_{01}e_{12})e_{23}$ to $e_{03}$.
The normalized tetrahedron family is $(L=R)$.

For a fixed horn, its paired fillers are the fibre of concatenation
by $\mathsf{ap}_{(-e_{23})}(q_{012})$. Concatenation by a path is an
equivalence, so this fibre is contractible. Thus normalization retains
the required horn condition without any truncation assumption on $E$.

Moreover, it is equivalent to the original $A_3$ family over each
boundary. To see this, eliminate the three given triangle paths. On
the canonical horn, the missing triangle selected by $A_3$ is exactly
the defining associator, and the normalized cell is its left unit path.
This gives a normalized cell at the original selected missing face
for every horn. Both families have contractible total spaces over the
same missing-face type and the same chosen base point. Applying the
fundamental theorem of identity types to both gives their fibrewise
equivalence. These steps are checked in \code{SemiSegalNormalize}.

\section{The pentagon is a missing tetrahedron}

The final input is a family $A_4$ over complete boundaries made from
these normalized tetrahedra. Require contractibility of
\[
 \sum_{t_{0234}} A_4(u,t_{0234})
\]
for each $\Lambda^4_1$ horn $u$. This is a condition on all matching
boundaries, not a separately specified associahedral equation.

For four composable edges $f,g,h,k$, form the horn with edges
\[
 e_{02}=fg,\quad e_{13}=gh,\quad e_{24}=hk,\quad
 e_{03}=f(gh),\quad e_{14}=g(hk),\quad e_{04}=f(g(hk)).
\]
Six triangles are reflexivity paths. The remaining four are
\begin{align*}
 q_{023}&=\alpha(f,g,h), & q_{134}&=\alpha(g,h,k),\\
 q_{024}&=\alpha(f,g,hk), &
 q_{034}&=\alpha(f,gh,k)\cdot\mathsf{ap}_{(f-)}\alpha(g,h,k).
\end{align*}
The four supplied tetrahedra are left unit paths. The missing
tetrahedron, selected by the 4-horn condition, says
\[
 \mathsf{ap}_{(-k)}\alpha(f,g,h)\cdot
   \bigl(\alpha(f,gh,k)\cdot\mathsf{ap}_{(f-)}\alpha(g,h,k)\bigr)
   =\alpha(fg,h,k)\cdot\alpha(f,g,hk).
\]
Reverse this equality and reassociate path concatenation to obtain
the exact \code{HcompPentagon.Pentagon} interface.
This is the checked theorem \code{SemiSegalPentagon.pentagon}.

This finite construction implements the low-dimensional semi-Segal
method, with normalization made explicit; compare Capriotti--Kraus,
\emph{Univalent Higher Categories via Complete Semi-Segal Types},
Section 3, \url{https://arxiv.org/abs/1707.03693}.
The simplicial boundary organization is compatible with the alternative
ambient presentation discussed by Sattler--W\"arn, but no ambient
degree-5 construction has been implemented here.

\section{Units and the homology interface}

Identity edges and two degeneracy triangles are additional inputs:
$A_2(1,f,f)$ and $A_2(f,1,f)$. They give the two unit laws by the
triangle equivalence. This is the unit fragment used by the homology
proofs. We do not assert a full degeneracy structure, completeness,
or Rezk univalence for this experimental presentation.

The names \code{Composition}, \code{CompLaws}, \code{Associativity},
and \code{HcompPentagon} expose the interface expected by the existing
theorems. A small inverse interface supplies the generic implication
from two-sided inverse to invertibility of postcomposition. In
\code{Homology.WithPentagon}, \code{hpent} is a definition using the
new pentagon theorem, not a segment parameter.

The existing theorem sources are reused by relative links. They cover
copowers, tensor maps, homotopy invariance, dimension, additivity,
Fubini, excision, groups, the positive-degree Mayer--Vietoris sequence,
reduced suspension, and the existing sphere calculations.
The copower and stability assumptions retain their universal-property
formulations. In particular, for a pushout of types, excision and
stability supply the pullback from which the sequence is constructed.

All three repeating positive-degree exactness positions use the same
structural boundary and group structures within this interpretation.
Degree-zero compatibility and categorical naturality remain open, and
no exactness is asserted at the final $H_1$ term. Sphere comparisons
retain their existing scope as bijections of sets.

\section{What has and has not been replaced}

The original 169 \code{hcat} source files are unchanged. In that root,
\code{hpent} is still assumed. In the experimental root, its replacement
is derived for the new operations and the homology proofs recheck.
No equality between the old and new group structures or connecting
maps is claimed: that would require comparison data relating the two
categorical presentations.

The experiment has explicit finite semi-Segal parameters, alongside
the unchanged foundational parameters \code{univ}, \code{images},
\code{surjections} and the homology universal-property inputs. No new
axioms have been installed in the original development. The index
type called \code{Cat} in the experiment merely indexes internal
presentations; it carries no ambient functor composition.
There are no judgmental associativity or unit rules for categories,
no UIP assumption, and no second equality.

The next task is to construct the finite internal data from an ambient
coherent category presentation, or to supply a model directly. The
refactor has removed the need to preserve the old associator in order
to rebuild homology. It has not removed the need for categorical
coherence itself.
\end{document}
